% Por que precisamos do geoide — alturas, números geopotenciais, IHRF e o Everest.
% Aula para o ensino médio, em "baby steps", com apêndice técnico.
%
% Compilar (nesta pasta):  latexmk -pdf geoide_e_alturas.tex
% Limpar auxiliares:       latexmk -c
%
% Acompanha o simulador historia_modelos_terrestres/ (abas Geóide e Teluróide).
% Segue o mesmo formato das aulas de ajusta_planos/explanations/.
%
% Fontes dos fatos e números citados (conferidos em 2026-09-21):
%   W0 = 62 636 853,4 m²/s²; coordenadas no ITRF; meta de ±1e-2 m²/s² (~1 mm):
%     IAG Resolution No. 1, IUGG General Assembly 2015 (Praga).
%     Sánchez et al., "Towards a first realization of the IHRS"; GGOS/IHRS Realization.
%   Grandes Lagos usam altitudes DINÂMICAS (IGLD 1985) justamente porque a água parada
%     tem altitudes ortométricas diferentes; NAVD 88 usa ortométricas de Helmert:
%     NOAA/NGS, "North American Vertical Datum and International Great Lakes Datum".
%   Everest: 8 848 m (Survey of India, 1954, neve); 8 844,43 m (China, 2005, rocha, com
%     radar de penetração no gelo); Nepal ao cume em 22/05/2019 (Khim Lal Gautam);
%     China em maio de 2020; anúncio conjunto em 08/12/2020 de 8 848,86 m, referido a um
%     geoide gravimétrico calculado com base no IHRS.
%   Ondulações extremas (-106 m no Índico, +85 m na Nova Guiné): mesmas do simulador,
%     models/geoide.js.

\documentclass[aspectratio=169,11pt]{beamer}

\usetheme[progressbar=frametitle, sectionpage=progressbar, numbering=fraction]{metropolis}
\metroset{block=fill}

\usepackage[T1]{fontenc}
\usepackage{lmodern}
\usepackage[brazilian]{babel}
\usepackage{icomma}
\usepackage{amsmath, amssymb, mathtools}
\usepackage{booktabs}
\usepackage[most]{tcolorbox}
\usepackage{tikz}
\usetikzlibrary{arrows.meta, calc, positioning, decorations.pathreplacing,
                decorations.pathmorphing, patterns, babel}

% ------------------------------------------------------------------ paleta (a mesma das outras aulas)
\definecolor{mStone}{HTML}{1C1917}
\definecolor{mStoneMid}{HTML}{78716C}
\definecolor{mStoneLight}{HTML}{F5F5F4}
\definecolor{mTeal}{HTML}{0D9488}
\definecolor{mRose}{HTML}{E11D48}
\definecolor{mAmber}{HTML}{D97706}
\definecolor{mIndigo}{HTML}{4F46E5}
\definecolor{mSky}{HTML}{0284C7}

% Cada superfície tem sua cor, a mesma em toda figura desta aula
\colorlet{corElip}{mIndigo}     % elipsoide — a forma lisa, o que o GNSS usa
\colorlet{corGeoide}{mTeal}     % geoide — a superfície que a água escolhe
\colorlet{corTerreno}{mAmber}   % terreno — a Terra de verdade
\colorlet{corAgua}{mSky}        % a água, e tudo que ela decide

\setbeamercolor{normal text}{fg=mStone, bg=white}
\setbeamercolor{palette primary}{fg=white, bg=mStone}
\setbeamercolor{frametitle}{fg=white, bg=mStone}
\setbeamercolor{alerted text}{fg=mRose}
\setbeamercolor{example text}{fg=mTeal}
\setbeamercolor{progress bar}{fg=mTeal, bg=mTeal!20}
\setbeamercolor{progress bar in head/foot}{fg=mTeal, bg=mTeal!20}
\setbeamercolor{progress bar in section page}{fg=mTeal, bg=mTeal!20}
\setbeamercolor{title separator}{fg=mTeal}
\setbeamercolor{block title}{fg=white, bg=mTeal}
\setbeamercolor{block body}{fg=mStone, bg=mTeal!8}

% ------------------------------------------------------------------ caixas
\newtcolorbox{ideia}[1][]{enhanced, colback=mTeal!7, colframe=mTeal, boxrule=0.6pt, arc=2pt,
  left=6pt, right=6pt, top=4pt, bottom=4pt, fonttitle=\bfseries\small, coltitle=white, #1}
\newtcolorbox{cuidado}[1][]{enhanced, colback=mAmber!8, colframe=mAmber, boxrule=0.6pt, arc=2pt,
  left=6pt, right=6pt, top=4pt, bottom=4pt, fonttitle=\bfseries\small, coltitle=white, #1}
\newtcolorbox{armadilha}[1][]{enhanced, colback=mRose!6, colframe=mRose, boxrule=0.6pt, arc=2pt,
  left=6pt, right=6pt, top=4pt, bottom=4pt, fonttitle=\bfseries\small, coltitle=white, #1}

% ------------------------------------------------------------------ estilos TikZ
\tikzset{
  seta/.style={-{Stealth[length=2.4mm]}, thick},
  rotulo/.style={font=\scriptsize, text=mStone},
  elip/.style={corElip, very thick},
  geoide/.style={corGeoide, very thick},
  terreno/.style={corTerreno, very thick},
  nivel/.style={corGeoide!55, thick},
}

\newcommand{\bnum}[1]{\textbf{\textcolor{mTeal}{#1}}}
\newcommand{\hgnss}{\textcolor{corElip}{h}}
\newcommand{\Horto}{\textcolor{corGeoide}{H}}
\newcommand{\Nond}{\textcolor{mRose}{N}}

\title{Por que a água às vezes corre para cima?}
\subtitle{O geoide, os sistemas de altitude, e a altura do Everest}
\author{}
\institute{Simulador de Modelos Terrestres}
\date{}

\begin{document}

\maketitle

% ==================================================================
\begin{frame}{O que vamos ver}
  \setbeamertemplate{section in toc}[sections numbered]
  \begin{columns}[T]
    \begin{column}{0.58\textwidth}
      \tableofcontents
    \end{column}
    \begin{column}{0.38\textwidth}
      \begin{ideia}[title={Duas partes}]
        \small
        \textbf{Parte 1} --- por que a altura precisa da gravidade, e o que é o geoide.
        \medskip

        \textbf{Parte 2} --- um zero para o mundo todo, e como isso resolveu a altura do
        Everest depois de décadas de briga.
      \end{ideia}
    \end{column}
  \end{columns}
\end{frame}

% ==================================================================
\section{Uma pergunta que parece boba}

\begin{frame}{``Qual é a altura daqui?''}
  \begin{columns}[T]
    \begin{column}{0.50\textwidth}
      Parece a pergunta mais simples do mundo. Mas experimente responder para um mesmo
      lugar, com instrumentos diferentes:
      \bigskip

      \centering
      \small
      \renewcommand{\arraystretch}{1.3}
      \begin{tabular}{lr}
        \toprule
        o celular (GPS) diz & $812{,}0$\,m \\
        a placa na estrada diz & $787{,}5$\,m \\
        o mapa antigo diz & $788{,}1$\,m \\
        \bottomrule
      \end{tabular}
    \end{column}
    \begin{column}{0.46\textwidth}
      \begin{cuidado}[title={Nenhum deles está errado}]
        Os três mediram bem. Eles só estão contando a altura \alert{a partir de lugares
        diferentes}.
      \end{cuidado}
      \bigskip

      Altura é sempre \bnum{altura a partir de alguma coisa}. Enquanto ninguém disser
      \emph{a partir de quê}, a pergunta não tem resposta.
    \end{column}
  \end{columns}
\end{frame}

\begin{frame}{Precisamos combinar onde fica o zero}
  \begin{columns}[c]
    \begin{column}{0.48\textwidth}
      Todo mundo já combinou isso uma vez: ``altura acima do \bnum{nível do mar}''.
      \bigskip

      Só que o mar \alert{não é plano}. Ele sobe, desce, tem maré, tem vento, tem água
      mais quente de um lado do que do outro.
      \bigskip

      E onde não há mar --- no meio de Goiás, no Himalaia --- \alert{onde fica o nível do
      mar?}
    \end{column}
    \begin{column}{0.48\textwidth}
      \centering
      \begin{tikzpicture}[scale=0.95]
        \fill[corAgua!18] (0,0) rectangle (6.2,-0.9);
        \draw[corAgua, thick, decorate,
              decoration={snake, amplitude=0.6mm, segment length=3mm}] (0,0) -- (6.2,0);
        \fill[corTerreno!25] (3.4,-0.9) .. controls (4.2,1.6) and (5.0,0.4) ..
                             (6.2,1.9) -- (6.2,-0.9) -- cycle;
        \draw[terreno] (3.4,-0.9) .. controls (4.2,1.6) and (5.0,0.4) .. (6.2,1.9);
        \node[rotulo, text=corAgua] at (1.4,0.42) {o mar, que não fica quieto};
        \node[rotulo, align=center] at (5.3,-0.45) {e aqui dentro?};
        \draw[mRose, seta] (5.3,-0.18) -- (5.3,0.9);
        \node[rotulo, text=mRose, right] at (5.4,1.2) {?};
      \end{tikzpicture}
    \end{column}
  \end{columns}
\end{frame}

\begin{frame}{O cano que não funciona}
  \begin{columns}[T]
    \begin{column}{0.42\textwidth}
      Uma obra precisa levar água de \textbf{A} até \textbf{B}, a 30\,km de distância.
      \medskip

      O GPS mede os dois pontos e dá a \alert{mesma altura}:
      \[ h_A = h_B = 812{,}00\ \text{m} \]
      Então a água deveria ficar parada.
      \medskip

      \onslide<2->{Só que na obra \alert{a água corre de A para B}.}
    \end{column}
    \begin{column}{0.54\textwidth}
      \centering
      \begin{tikzpicture}[scale=0.86]
        \draw[corElip, dashed, thick] (0,0) -- (6.4,0);
        \node[rotulo, text=corElip, right] at (5.0,-0.32) {mesma altura pelo GPS};
        \fill[mStone] (0.5,0) circle (0.08) node[above, rotulo] {A};
        \fill[mStone] (5.9,0) circle (0.08) node[above, rotulo] {B};
        \onslide<2->{
          \draw[corAgua, line width=2pt, -{Stealth[length=3mm]}] (0.9,0.28) -- (5.5,0.28);
          \node[rotulo, text=corAgua, above] at (3.2,0.42) {a água vai para lá};
          \draw[corGeoide, very thick] (0,-1.5) .. controls (2.2,-1.2) and (4.2,-1.9) ..
                                       (6.4,-1.75);
          \node[rotulo, text=corGeoide, below] at (3.2,-1.95) {o ``zero'' de verdade não é plano};
          \draw[corGeoide!60] (0.5,0) -- (0.5,-1.42);
          \draw[corGeoide!60] (5.9,0) -- (5.9,-1.78);
        }
      \end{tikzpicture}
    \end{column}
  \end{columns}
  \onslide<2->{
    \medskip
    \begin{armadilha}[title={O que isso significa}]
      A altura que o GPS dá é uma \alert{distância geométrica}, medida a partir de uma
      forma lisa inventada. A água não obedece a formas inventadas: ela obedece à
      \bnum{gravidade}.
    \end{armadilha}
  }
\end{frame}

% ==================================================================
\section{Três superfícies, três alturas}

\begin{frame}{A Terra de verdade é cheia de calombos}
  \begin{columns}[c]
    \begin{column}{0.48\textwidth}
      A superfície onde a gente pisa --- o \textcolor{corTerreno}{\textbf{terreno}} ---
      tem montanhas de 8\,km e fossas de 11\,km.
      \bigskip

      Ninguém consegue escrever uma fórmula para ela. Então a geodésia faz o que sempre
      se faz em ciência: \bnum{inventa formas simples que chegam perto}, uma de cada vez.
    \end{column}
    \begin{column}{0.48\textwidth}
      \centering
      \begin{tikzpicture}[scale=0.92]
        \draw[terreno] (0,0) .. controls (0.8,0.9) and (1.4,-0.4) .. (2.2,0.35)
                       .. controls (2.9,1.0) and (3.4,0.1) .. (4.0,1.5)
                       .. controls (4.6,0.2) and (5.4,0.7) .. (6.2,0.1);
        \fill[corTerreno!18] (0,0) .. controls (0.8,0.9) and (1.4,-0.4) .. (2.2,0.35)
                       .. controls (2.9,1.0) and (3.4,0.1) .. (4.0,1.5)
                       .. controls (4.6,0.2) and (5.4,0.7) .. (6.2,0.1)
                       -- (6.2,-0.8) -- (0,-0.8) -- cycle;
        \node[rotulo, text=corTerreno, above] at (4.0,1.55) {terreno};
      \end{tikzpicture}
    \end{column}
  \end{columns}
\end{frame}

\begin{frame}{Forma nº 1: o elipsoide --- a Terra lisa de mentirinha}
  \begin{columns}[T]
    \begin{column}{0.50\textwidth}
      Uma bola levemente achatada nos polos, descrita por \bnum{dois números só}:
      \[ a = 6\,378\,137\ \text{m}, \qquad f = 1/298{,}257 \]
      \medskip
      É \alert{pura geometria}: nenhuma pedra, nenhuma montanha, nenhuma gravidade.
      Uma fórmula, e pronto.
      \medskip

      É esta a forma que o \bnum{GPS} usa. A altura que o celular mostra é a distância
      até ela: a \textcolor{corElip}{\textbf{altitude elipsoidal $h$}}.
    \end{column}
    \begin{column}{0.46\textwidth}
      \centering
      \begin{tikzpicture}[scale=0.95]
        \draw[elip] (3,0) ellipse (2.6 and 2.1);
        \draw[mStoneMid, dashed] (0.4,0) -- (5.6,0);
        \draw[mStoneMid, dashed] (3,-2.1) -- (3,2.1);
        \draw[corElip, {Stealth[length=2mm]}-{Stealth[length=2mm]}] (3,0) -- (5.6,0);
        \node[rotulo, text=corElip, above] at (4.4,0.05) {$a$};
        \node[rotulo, text=corElip, right] at (3.05,1.15) {$b$};
        \draw[corElip, {Stealth[length=2mm]}-{Stealth[length=2mm]}] (3,0) -- (3,2.1);
        \node[rotulo, below] at (3,-2.35) {achatado nos polos (bem pouco)};
      \end{tikzpicture}
    \end{column}
  \end{columns}
\end{frame}

\begin{frame}{O GPS não sabe onde é ``para baixo''}
  \begin{columns}[c]
    \begin{column}{0.52\textwidth}
      Um receptor GNSS resolve um problema puramente geométrico: \emph{onde estou, em
      relação a um punhado de satélites?}
      \medskip

      Ele devolve latitude, longitude e $h$ --- e \alert{nunca mediu a gravidade}.
      Não tem como: não há nada dentro dele que sinta peso.
      \medskip

      \begin{cuidado}[title={A consequência}]
        $h$ diz a que \emph{distância} você está do elipsoide. Não diz se a água vai subir
        ou descer indo daqui até ali.
      \end{cuidado}
    \end{column}
    \begin{column}{0.44\textwidth}
      \centering
      \begin{tikzpicture}[scale=0.9]
        \foreach \x/\y in {0.4/2.9, 2.2/3.4, 4.2/3.1, 5.6/2.4} {
          \fill[mStoneMid] (\x,\y) circle (0.11);
          \draw[mStoneMid!60] (\x,\y) -- (3.0,0.55);
        }
        \fill[mStone] (3.0,0.5) rectangle (3.3,0.75);
        \draw[corElip, very thick] (0.2,0) .. controls (2,0.25) and (4,0.25) .. (5.9,0);
        \draw[corElip, {Stealth[length=2mm]}-{Stealth[length=2mm]}] (3.15,0.2) -- (3.15,0.5);
        \node[rotulo, text=corElip, right] at (3.25,0.35) {$h$};
        \node[rotulo, text=corElip, below] at (3.0,-0.2) {elipsoide};
        \node[rotulo] at (2.0,2.2) {satélites};
      \end{tikzpicture}
    \end{column}
  \end{columns}
\end{frame}

\begin{frame}{Forma nº 2: o que a água faz quando é deixada em paz}
  \begin{columns}[T]
    \begin{column}{0.50\textwidth}
      Encha uma bacia e espere. A água para numa superfície bem definida ---
      a que \bnum{não tem nenhuma ladeira}.
      \bigskip

      Agora imagine isso no planeta inteiro: tire o vento, a maré, as correntes, e deixe o
      oceano se acalmar. A superfície que sobra é uma superfície especial.
      \bigskip

      E ela \bnum{continua por baixo dos continentes}: é onde a água pararia se houvesse
      canaletas cortando a terra toda.
    \end{column}
    \begin{column}{0.46\textwidth}
      \centering
      \begin{tikzpicture}[scale=0.95]
        \fill[corTerreno!25] (0,0) .. controls (1.2,0.15) and (1.8,1.4) .. (3.1,1.45)
              .. controls (4.4,1.5) and (5.0,0.2) .. (6.2,0.1) -- (6.2,-1.2) -- (0,-1.2) -- cycle;
        \draw[terreno] (0,0) .. controls (1.2,0.15) and (1.8,1.4) .. (3.1,1.45)
              .. controls (4.4,1.5) and (5.0,0.2) .. (6.2,0.1);
        \draw[corAgua, dashed, very thick] (0,-0.25) -- (6.2,-0.25);
        \fill[corAgua!35] (2.55,-0.25) rectangle (2.95,1.46);
        \fill[corAgua!35] (0,-0.25) rectangle (0.6,-0.05);
        \node[rotulo, text=corAgua, below right] at (0.1,-0.4) {a mesma superfície, por baixo de tudo};
        \node[rotulo, text=corTerreno, above] at (3.1,1.5) {terreno};
      \end{tikzpicture}
    \end{column}
  \end{columns}
\end{frame}

\begin{frame}{Essa superfície tem nome: geoide}
  \begin{columns}[T]
    \begin{column}{0.52\textwidth}
      \begin{ideia}[title={Geoide}]
        A superfície em que a gravidade tem sempre \bnum{o mesmo valor de potencial} ---
        aquela que coincide com o nível médio dos mares em repouso, e que segue por baixo
        dos continentes.
      \end{ideia}
      \medskip

      O nome foi cunhado por \textbf{J. B. Listing} em 1873: do grego \emph{ge} (Terra)
      $+$ \emph{eidos} (forma). Literalmente: \emph{``a forma da Terra''}.
      \medskip

      Ele é \alert{ondulado}: onde há mais massa, o geoide sobe; onde há menos, afunda.
    \end{column}
    \begin{column}{0.44\textwidth}
      \centering
      \begin{tikzpicture}[scale=0.92]
        \draw[elip, dashed] (3,0) ellipse (2.15 and 1.80);
        \draw[geoide] plot[smooth cycle, tension=0.8]
          coordinates {(5.35,0.55) (4.45,1.55) (2.95,2.35) (1.35,1.35) (0.45,0.05)
                       (1.05,-1.30) (2.55,-2.15) (4.35,-1.20)};
        \node[rotulo, text=corElip, right] at (4.85,-1.75) {elipsoide};
        \node[rotulo, text=corGeoide, above] at (2.95,2.40) {geoide};
      \end{tikzpicture}
      \smallskip

      {\scriptsize\textcolor{mStoneMid}{As ondulações estão \emph{muito} exageradas:\\
      de verdade elas são milhares de vezes menores.}}
    \end{column}
  \end{columns}
\end{frame}

\begin{frame}{Quanto o geoide realmente ondula}
  \begin{columns}[T]
    \begin{column}{0.50\textwidth}
      A distância entre o geoide e o elipsoide é a \textcolor{mRose}{\textbf{ondulação
      geoidal $N$}}. No planeta inteiro ela vai de:
      \medskip
      \centering
      \small
      \renewcommand{\arraystretch}{1.25}
      \begin{tabular}{lr}
        \toprule
        Depressão do Oceano Índico & $-106$\,m \\
        Domo da Nova Guiné & $+85$\,m \\
        Alto da Islândia & $+68$\,m \\
        Baía de Hudson & $-58$\,m \\
        \bottomrule
      \end{tabular}
    \end{column}
    \begin{column}{0.46\textwidth}
      \begin{cuidado}[title={Parece muito? É pouquíssimo.}]
        $106$ metros num raio de $6\,371$\,\textbf{quilômetros}.
        \medskip

        É como um calombo de \bnum{1 milímetro} numa bola de basquete --- e mesmo assim
        ninguém pode ignorá-lo numa obra.
      \end{cuidado}
      \medskip

      {\small É por isso que toda figura de geoide que você já viu está com o relevo
      exagerado umas 10\,000 vezes.}
    \end{column}
  \end{columns}
\end{frame}

\begin{frame}{A equação mais importante desta aula}
  \begin{columns}[T]
    \begin{column}{0.42\textwidth}
      Três superfícies, empilhadas:
      \begin{itemize}
        \item \textcolor{corTerreno}{\textbf{terreno}} --- onde você está
        \item \textcolor{corGeoide}{\textbf{geoide}} --- o zero físico
        \item \textcolor{corElip}{\textbf{elipsoide}} --- o zero geométrico
      \end{itemize}
      \medskip
      \[ \boxed{\; \hgnss = \Horto + \Nond \;} \]
      \begin{itemize}
        \item $\hgnss$: o que o GPS dá
        \item $\Horto$: a altitude que a placa mostra
        \item $\Nond$: a ondulação, de um modelo
      \end{itemize}
    \end{column}
    \begin{column}{0.54\textwidth}
      \centering
      \begin{tikzpicture}[scale=1.0]
        \draw[terreno] (0,2.0) .. controls (1.6,2.5) and (3.2,1.6) .. (6.4,2.2);
        \node[rotulo, text=corTerreno, above left] at (6.4,2.25) {terreno};
        \draw[geoide] (0,0.75) .. controls (2.2,1.1) and (4.2,0.4) .. (6.4,0.7);
        \node[rotulo, text=corGeoide, below left] at (6.4,0.6) {geoide};
        \draw[elip] (0,-0.5) -- (6.4,-0.5);
        \node[rotulo, text=corElip, above left] at (6.4,-0.45) {elipsoide};
        % o ponto e as tres medidas
        \fill[mStone] (2.6,2.27) circle (0.09);
        \node[rotulo, above] at (2.6,2.36) {você está aqui};
        \draw[corGeoide, {Stealth[length=2mm]}-{Stealth[length=2mm]}] (2.45,2.27) -- (2.45,0.93);
        \node[rotulo, text=corGeoide, left] at (2.38,1.6) {$H$};
        \draw[mRose, {Stealth[length=2mm]}-{Stealth[length=2mm]}] (2.9,0.93) -- (2.9,-0.5);
        \node[rotulo, text=mRose, right] at (2.97,0.2) {$N$};
        \draw[corElip, {Stealth[length=2mm]}-{Stealth[length=2mm]}] (3.6,2.13) -- (3.6,-0.5);
        \node[rotulo, text=corElip, right] at (3.67,0.85) {$h$};
        \draw[mStoneMid, dotted] (2.6,2.27) -- (3.6,2.13);
      \end{tikzpicture}
    \end{column}
  \end{columns}
\end{frame}

\begin{frame}{Voltando ao cano que não funcionava}
  \centering
  \small
  \renewcommand{\arraystretch}{1.35}
  \begin{tabular}{lrrr}
    \toprule
    & \textcolor{corElip}{$h$ (GPS)} & \textcolor{mRose}{$N$ (modelo)}
    & \textcolor{corGeoide}{$H = h - N$} \\
    \midrule
    ponto A & $812{,}00$\,m & $-8{,}40$\,m & \textcolor{corGeoide}{$\mathbf{820{,}40}$\,m} \\
    ponto B & $812{,}00$\,m & $-8{,}12$\,m & \textcolor{corGeoide}{$\mathbf{820{,}12}$\,m} \\
    \midrule
    diferença & $0{,}00$\,m & & \alert{$\mathbf{0{,}28}$\,m} \\
    \bottomrule
  \end{tabular}
  \bigskip

  A água desce \bnum{28 centímetros} de A para B --- exatamente o que se via na obra.
  \medskip

  \begin{ideia}[title={O que o geoide resolveu}]
    O GPS estava certo, e a água também. Faltava a \bnum{tradução} entre os dois:
    a ondulação $N$.
  \end{ideia}
\end{frame}

% ==================================================================
\section{O que quer dizer ``estar no mesmo nível''}

\begin{frame}{Superfície de nível: onde a água não corre}
  \begin{columns}[T]
    \begin{column}{0.50\textwidth}
      O geoide não é uma superfície especial e solitária. Ele é \bnum{uma de uma família
      inteira}.
      \medskip

      Acima e abaixo dele existem outras superfícies com a mesma propriedade: soltar uma
      bolinha ali, e ela não rola. São as \bnum{superfícies de nível}.
      \medskip

      O geoide é simplesmente aquela que escolhemos chamar de \alert{zero}.
    \end{column}
    \begin{column}{0.46\textwidth}
      \centering
      \begin{tikzpicture}[scale=0.82]
        \foreach \y/\amp in {0/0.30, 0.75/0.27, 1.5/0.24, 2.25/0.21, 3.0/0.18} {
          \draw[nivel] (0,\y) .. controls (2.1,\y+\amp) and (4.2,\y-\amp) .. (6.3,\y+0.5*\amp);
        }
        \draw[geoide] (0,0) .. controls (2.1,0.30) and (4.2,-0.30) .. (6.3,0.15);
        \node[rotulo, text=corGeoide, below] at (5.2,-0.15) {geoide};
        \node[rotulo, align=left, text=mStoneMid] at (2.0,3.7) {superfícies de nível};
      \end{tikzpicture}
    \end{column}
  \end{columns}
\end{frame}

\begin{frame}{O fio de prumo aponta para o geoide, não para o elipsoide}
  \begin{columns}[T]
    \begin{column}{0.50\textwidth}
      Um fio de prumo fica perpendicular à superfície de nível que passa por ele ---
      ou seja, \bnum{perpendicular ao geoide}.
      \medskip

      A perpendicular ao \emph{elipsoide} é outra direção. O ângulo entre as duas é o
      \bnum{desvio da vertical}.
      \medskip

      Em terreno plano ele é de poucos segundos de arco. Perto de uma cordilheira,
      \alert{a montanha puxa o prumo para si} e o desvio fica bem maior.
    \end{column}
    \begin{column}{0.46\textwidth}
      \centering
      \begin{tikzpicture}[scale=0.84]
        \draw[geoide] (0,1.2) .. controls (2.0,1.75) and (4.0,0.9) .. (6.0,1.35);
        \node[rotulo, text=corGeoide, right] at (5.2,0.98) {geoide};
        \draw[elip] (0,0) -- (6.0,0);
        \node[rotulo, text=corElip, right] at (5.2,-0.25) {elipsoide};
        \fill[mStone] (2.6,1.62) circle (0.08);
        % normal ao elipsoide
        \draw[corElip, thick, dashed] (2.6,1.62) -- (2.6,3.2);
        % normal ao geoide (inclinada)
        \draw[corGeoide, very thick] (2.6,1.62) -- (3.05,3.15);
        \draw[mRose, thick] (2.6,2.45) arc (90:73:0.85);
        \node[rotulo, text=mRose, left] at (2.05,2.28) {desvio};
        \draw[mRose, seta] (2.15,2.38) -- (2.58,2.55);
        \node[rotulo, text=corElip, above, align=center] at (1.95,3.15)
              {normal ao\\elipsoide};
        \draw[corElip, seta, dashed] (2.20,3.20) -- (2.56,3.20);
        \node[rotulo, text=corGeoide, right] at (3.12,3.15) {prumo};
      \end{tikzpicture}
    \end{column}
  \end{columns}
\end{frame}

\begin{frame}{E aqui vem a surpresa: elas não são paralelas}
  \begin{columns}[T]
    \begin{column}{0.48\textwidth}
      Duas superfícies de nível vizinhas \alert{não guardam a mesma distância} uma da
      outra em todo lugar.
      \medskip

      Onde a gravidade é mais forte --- nos polos --- elas ficam \bnum{mais juntas}.
      Onde é mais fraca --- no Equador --- ficam \bnum{mais afastadas}.
      \medskip

      {\small A gravidade varia cerca de $0{,}5\,\%$ do Equador ao polo: o mesmo saco de
      cimento pesa meio por cento a mais lá em cima.}
    \end{column}
    \begin{column}{0.48\textwidth}
      \centering
      \begin{tikzpicture}[scale=0.80]
        \draw[geoide] (3,0) ellipse (2.5 and 2.05);
        \draw[nivel] (3,0) ellipse (2.92 and 2.30);
        \draw[nivel] (3,0) ellipse (3.32 and 2.53);
        \draw[mRose, {Stealth[length=1.8mm]}-{Stealth[length=1.8mm]}] (3,2.05) -- (3,2.30);
        \draw[mRose, {Stealth[length=1.8mm]}-{Stealth[length=1.8mm]}] (5.5,0) -- (5.92,0);
        \node[rotulo, text=mRose, above] at (3,2.60) {mais juntas};
        \node[rotulo, text=mRose, right, align=left] at (6.15,0) {mais\\afastadas};
        \node[rotulo] at (3,-2.45) {(exagerado)};
      \end{tikzpicture}
    \end{column}
  \end{columns}
\end{frame}

\begin{frame}{Por que isso estraga a ideia ingênua de altura}
  Se as superfícies de nível não são paralelas, então \alert{``subir 1 metro'' não quer
  dizer sempre a mesma coisa}.
  \bigskip
  \begin{columns}[T]
    \begin{column}{0.48\textwidth}
      \begin{armadilha}[title={O nivelamento que não fecha}]
        Meça a diferença de altura entre dois pontos indo por um caminho, e depois por
        outro. \alert{Dá diferente} --- mesmo com instrumentos perfeitos.
        \medskip

        Não é erro de medida: é o mundo sendo assim.
      \end{armadilha}
    \end{column}
    \begin{column}{0.48\textwidth}
      \begin{ideia}[title={A saída}]
        Parar de medir altura em \emph{metros subidos} e passar a medir
        \bnum{o esforço para subir}.
        \medskip

        Esse esforço não depende do caminho --- e é a próxima seção.
      \end{ideia}
    \end{column}
  \end{columns}
\end{frame}

% ==================================================================
\section{A altura medida em esforço}

\begin{frame}{Subir cansa --- e é isso que importa}
  \begin{columns}[T]
    \begin{column}{0.50\textwidth}
      Carregue um saco de cimento escada acima. O que cansa não é ``ter subido 3 metros'':
      é o \bnum{trabalho} que você fez contra a gravidade.
      \medskip

      Onde a gravidade é mais forte, os mesmos 3 metros \alert{cansam mais}.
      \medskip

      A geodésia moderna adota exatamente essa medida. Ela chama de
      \bnum{número geopotencial}.
    \end{column}
    \begin{column}{0.46\textwidth}
      \centering
      \begin{tikzpicture}[scale=0.95]
        \foreach \i in {0,...,3} {
          \draw[mStoneMid, thick] (\i*0.65,\i*0.5) -- (\i*0.65+0.65,\i*0.5)
                                  -- (\i*0.65+0.65,\i*0.5+0.5);
        }
        \fill[corTerreno] (2.05,1.62) circle (0.17);
        \draw[corTerreno, line width=1.6pt] (2.05,1.45) -- (2.05,0.95);
        \draw[corTerreno, thick] (2.05,1.30) -- (1.70,1.05) (2.05,1.30) -- (2.40,1.05);
        \draw[corTerreno, thick] (2.05,0.95) -- (1.80,0.5) (2.05,0.95) -- (2.35,0.5);
        \draw[mRose, seta, line width=1.4pt] (3.4,2.3) -- (3.4,0.55);
        \node[rotulo, text=mRose, right] at (3.5,1.5) {gravidade};
        \node[rotulo, align=center] at (1.5,2.6) {o que cansa é o \emph{trabalho},\\não a distância};
      \end{tikzpicture}
    \end{column}
  \end{columns}
\end{frame}

\begin{frame}{O número geopotencial}
  \begin{columns}[T]
    \begin{column}{0.52\textwidth}
      Dê um nome ao ``nível de energia'' de cada superfície de nível: chame de $W$.
      O do geoide é $W_0$.
      \medskip

      O \bnum{número geopotencial} de um ponto é a diferença:
      \[ \boxed{\; C = W_0 - W \;} \]
      \medskip

      Em palavras: \emph{o trabalho que é preciso fazer, por quilo carregado, para levar
      alguma coisa do geoide até ali} --- medido em $\text{m}^2/\text{s}^2$, que é
      energia por quilo, \alert{e não em metros}.
    \end{column}
    \begin{column}{0.44\textwidth}
      \begin{ideia}[title={As três propriedades que salvam tudo}]
        \begin{enumerate}
          \item $C$ é \bnum{único} para cada ponto;
          \item \alert{não depende do caminho} percorrido para chegar lá;
          \item dois pontos com o mesmo $C$ estão \bnum{mesmo} no mesmo nível ---
                a água entre eles fica parada.
        \end{enumerate}
      \end{ideia}
    \end{column}
  \end{columns}
\end{frame}

\begin{frame}{De ``esforço'' de volta para metros}
  Ninguém quer uma placa de estrada escrita em $\text{m}^2/\text{s}^2$. Então divide-se
  $C$ por uma gravidade --- e a escolha da gravidade \alert{dá nomes diferentes}:
  \medskip
  \centering
  \small
  \renewcommand{\arraystretch}{1.5}
  \begin{tabular}{p{3.6cm}lp{4.6cm}}
    \toprule
    dividindo $C$ por & resulta & onde se usa \\
    \midrule
    a gravidade \textbf{real} média & altitude
      \textcolor{corGeoide}{\textbf{ortométrica}} $H$ &
      a altitude ``acima do nível do mar'' clássica \\
    a gravidade \textbf{teórica} do elipsoide & altitude \textbf{normal} $H^{*}$ &
      adotada por muitos países, inclusive o Brasil \\
    uma gravidade \textbf{constante} & altitude \textbf{dinâmica} &
      quando o que importa é a água, não a distância \\
    \bottomrule
  \end{tabular}
  \medskip

  As três saem do \bnum{mesmo $C$}. Mudou só a régua usada para traduzi-lo.
\end{frame}

\begin{frame}{Um caso real: os Grandes Lagos}
  \begin{columns}[T]
    \begin{column}{0.52\textwidth}
      O Lago Superior é um só corpo de água parada. Mesmo assim, as pontas dele têm
      \alert{altitudes ortométricas diferentes}.
      \medskip

      Isso não é erro: é a consequência de as superfícies de nível não serem paralelas.
      Mas é um pesadelo para quem opera eclusas, hidrelétricas e cartas náuticas.
      \medskip

      Por isso o datum dos Grandes Lagos --- o \textbf{IGLD 1985} --- abandonou as
      altitudes ortométricas e usa \bnum{altitudes dinâmicas}.
    \end{column}
    \begin{column}{0.44\textwidth}
      \centering
      \begin{tikzpicture}[scale=0.95]
        \fill[corAgua!25] (0.2,0) rectangle (5.6,0.85);
        \draw[corAgua, very thick] (0.2,0.85) -- (5.6,0.85);
        \node[rotulo, text=corAgua] at (2.9,0.42) {água parada};
        \draw[corGeoide, dashed] (0.2,1.35) -- (5.6,1.62);
        \node[rotulo, text=corGeoide, above, align=center] at (2.9,1.55)
              {superfície de nível\\(inclinada em relação à água)};
        \node[rotulo, text=mRose, below] at (1.0,-0.1) {$H = 183{,}2$};
        \node[rotulo, text=mRose, below] at (4.8,-0.1) {$H = 183{,}0$};
        \node[rotulo, below] at (2.9,-0.55) {mesmo lago, ``alturas'' diferentes};
      \end{tikzpicture}
    \end{column}
  \end{columns}
\end{frame}

\begin{frame}{Como se conhece o geoide}
  \begin{columns}[T]
    \begin{column}{0.46\textwidth}
      \begin{itemize}
        \item \textbf{Gravímetros em campo} --- medem $g$ ponto a ponto, no chão, em
              navios, em aviões.
        \item \textbf{Satélites dedicados} --- \textbf{GRACE} (NASA/DLR) mede a distância
              entre dois satélites gêmeos, que se afasta e se aproxima conforme a massa lá
              embaixo; \textbf{GOCE} (ESA) levou um gradiômetro em órbita baixa.
        \item \textbf{Modelos globais} --- o \textbf{EGM2008} junta tudo numa fórmula com
              milhões de termos.
      \end{itemize}
    \end{column}
    \begin{column}{0.50\textwidth}
      \centering
      \begin{tikzpicture}[scale=0.92]
        \draw[geoide] (3,0) ellipse (1.75 and 1.45);
        \draw[mStoneMid, dashed] (3,0) ellipse (2.85 and 2.35);
        \fill[mStone] (1.05,1.05) rectangle (1.45,1.25);
        \fill[mStone] (1.75,1.55) rectangle (2.15,1.75);
        \draw[mRose, {Stealth[length=1.6mm]}-{Stealth[length=1.6mm]}] (1.5,1.16) -- (1.72,1.5);
        \node[rotulo, right, align=left] at (2.45,2.25)
              {GRACE: dois satélites\\que ``se sentem''};
        \draw[corTerreno, seta] (4.75,-2.15) -- (4.18,-1.30);
        \node[rotulo, text=corTerreno, right] at (4.6,-2.3) {gravímetro no chão};
      \end{tikzpicture}
    \end{column}
  \end{columns}
\end{frame}

\begin{frame}[standout]
  Parte 2\\[6pt]
  {\normalsize Um zero só para o mundo inteiro}
\end{frame}

% ==================================================================
\section{Cada país com o seu próprio zero}

\begin{frame}{O zoológico dos zeros}
  \begin{columns}[T]
    \begin{column}{0.48\textwidth}
      Cada país escolheu, lá atrás, \bnum{um marégrafo} num porto seu, mediu o nível médio
      do mar por alguns anos e cravou: \emph{``aqui é o zero''}.
      \medskip

      Todos fizeram a coisa certa. Mas:
      \begin{itemize}
        \item cada mar está num nível um pouco diferente;
        \item cada país mediu numa época diferente;
        \item o mar subiu desde então.
      \end{itemize}
      \medskip

      Resultado: os zeros do mundo diferem entre si por até \alert{uns 2 metros}.
    \end{column}
    \begin{column}{0.48\textwidth}
      \centering
      \begin{tikzpicture}[scale=0.95]
        \fill[corTerreno!18] (0,0.62) rectangle (6.3,2.0);
        \draw[terreno] (0,2.0) -- (6.3,2.0);
        \foreach \x/\y/\lab in {0.78/0.44/A, 2.36/0.12/B, 3.94/0.52/C, 5.52/0.24/D} {
          \draw[corAgua, very thick] (\x-0.6,\y) -- (\x+0.6,\y);
          \draw[mStoneMid, thin] (\x,\y) -- (\x,2.0);
          \fill[mStone] (\x-0.09,2.0) rectangle (\x+0.09,2.22);
          \node[rotulo, below] at (\x,-0.14) {país \lab};
        }
        \draw[mStoneMid, dashed] (0,0.32) -- (6.3,0.32);
        \node[rotulo, align=center] at (3.15,2.75) {cada marégrafo crava um zero diferente};
      \end{tikzpicture}
    \end{column}
  \end{columns}
\end{frame}

\begin{frame}{A ponte que não encaixa}
  \begin{columns}[c]
    \begin{column}{0.50\textwidth}
      Enquanto cada país trabalhava sozinho, isso não incomodava ninguém.
      \medskip

      Hoje incomoda muito:
      \begin{itemize}
        \item obras que cruzam fronteiras --- pontes, túneis, dutos;
        \item bacias hidrográficas compartilhadas;
        \item medir a subida do nível do mar \alert{no planeta}, e não em cada porto;
        \item o GNSS dá $h$ em qualquer lugar do mundo --- e aí é preciso um $N$ que
              também valha no mundo todo.
      \end{itemize}
    \end{column}
    \begin{column}{0.46\textwidth}
      \centering
      \begin{tikzpicture}[scale=0.95]
        \fill[corTerreno!22] (0,0) rectangle (2.1,1.3);
        \fill[corTerreno!22] (4.2,0) rectangle (6.3,1.3);
        \draw[terreno] (0,1.3) -- (2.1,1.3);
        \draw[terreno] (4.2,1.55) -- (6.3,1.55);
        \draw[mStone, line width=2pt] (2.1,1.3) -- (4.2,1.3);
        \draw[mRose, line width=2pt] (4.2,1.55) -- (4.2,1.3);
        \node[rotulo, text=mRose, above, align=center] at (3.1,1.95)
              {degrau de $25$\,cm};
        \draw[mRose, seta] (3.6,1.92) -- (4.14,1.48);
        \node[rotulo, below] at (3.15,0.9) {rio de fronteira};
      \end{tikzpicture}
    \end{column}
  \end{columns}
\end{frame}

% ==================================================================
\section{IHRS e IHRF: o zero combinado}

\begin{frame}{A ideia: em vez de um porto, um número}
  \begin{columns}[T]
    \begin{column}{0.52\textwidth}
      Escolher um marégrafo sempre vai favorecer alguém e deixar o resto do mundo torto.
      \medskip

      Então a comunidade geodésica fez outra coisa: escolheu \bnum{um valor de potencial}
      --- um número puro, sem país, sem porto, sem data.
      \medskip

      Em 2015, na Assembleia Geral da IUGG em Praga, a \textbf{IAG} aprovou a
      \textbf{Resolução nº 1}, criando o \bnum{IHRS}: o Sistema Internacional de
      Referência de Altitudes.
    \end{column}
    \begin{column}{0.44\textwidth}
      \begin{ideia}[title={O zero do mundo}]
        \centering
        \[ W_0 = 62\,636\,853{,}4\ \frac{\text{m}^2}{\text{s}^2} \]
      \end{ideia}
      \medskip

      {\small A superfície de nível que tem exatamente esse potencial passa a ser
      \alert{o geoide oficial do planeta}. Ela não pertence a ninguém --- e serve a todos.}
    \end{column}
  \end{columns}
\end{frame}

\begin{frame}{As peças do IHRS}
  \centering
  \small
  \renewcommand{\arraystretch}{1.45}
  \begin{tabular}{lp{7.4cm}}
    \toprule
    \textbf{o zero} & a superfície de nível com $W_0 = 62\,636\,853{,}4\ \text{m}^2/\text{s}^2$ \\
    \textbf{a altura} & o \bnum{número geopotencial} $C_P = W_0 - W(P)$ --- e não metros \\
    \textbf{a posição} & coordenadas no \textbf{ITRF}, o mesmo referencial do GNSS \\
    \textbf{a meta} & $\pm 0{,}01\ \text{m}^2/\text{s}^2$, que é cerca de
      \bnum{$\pm 1$ milímetro} de altura \\
    \bottomrule
  \end{tabular}
  \bigskip

  \begin{ideia}[title={Repare na escolha}]
    A altura oficial do mundo é medida em \bnum{esforço}, não em metros --- exatamente
    pelo motivo da Parte 1: só o esforço não depende do caminho.
  \end{ideia}
\end{frame}

\begin{frame}{Sistema e materialização: IHRS e IHRF}
  \begin{columns}[T]
    \begin{column}{0.48\textwidth}
      \begin{ideia}[title={IHRS --- o sistema}]
        A \bnum{regra}. Diz o que é o zero, o que é ``altura'' e em que referencial as
        posições são dadas.
        \medskip

        É um texto. Não dá para medir nada com ele.
      \end{ideia}
    \end{column}
    \begin{column}{0.48\textwidth}
      \begin{ideia}[title={IHRF --- o frame}]
        A \bnum{rede de estações} de verdade, espalhadas pelo planeta, com número
        geopotencial conhecido.
        \medskip

        É onde se amarra o trabalho de cada país.
      \end{ideia}
    \end{column}
  \end{columns}
  \bigskip
  {\small A mesma dupla existe do lado horizontal: o \emph{sistema} diz como as
  coordenadas são definidas; o \emph{frame} (o ITRF) é a lista de estações com
  coordenadas medidas. O IHRF segue a mesma estrutura: uma rede global, densificada por
  redes regionais e nacionais.}
\end{frame}

% ==================================================================
\section{O Everest, enfim resolvido}

\begin{frame}{Duas respostas para a mesma montanha}
  \begin{columns}[T]
    \begin{column}{0.50\textwidth}
      Por décadas, a montanha mais alta do mundo teve \alert{duas alturas oficiais} ---
      e cada país insistia na sua.
      \medskip
      \centering
      \small
      \renewcommand{\arraystretch}{1.35}
      \begin{tabular}{lr}
        \toprule
        Nepal (Survey of India, 1954) & $8\,848$\,m \\
        China (2005) & $8\,844{,}43$\,m \\
        \midrule
        diferença & \alert{$3{,}57$\,m} \\
        \bottomrule
      \end{tabular}
    \end{column}
    \begin{column}{0.46\textwidth}
      \begin{cuidado}[title={Não era teimosia}]
        Os dois países mediram com cuidado e competência.
        \medskip

        Eles discordavam porque estavam respondendo a \bnum{duas perguntas diferentes} ---
        e contando a partir de \bnum{dois zeros diferentes}.
      \end{cuidado}
    \end{column}
  \end{columns}
\end{frame}

\begin{frame}{Primeira discordância: até onde é a montanha?}
  \begin{columns}[T]
    \begin{column}{0.48\textwidth}
      O cume do Everest é coberto por uma \bnum{calota de neve e gelo} de alguns metros.
      \medskip

      \begin{itemize}
        \item A China mediu até a \textcolor{corTerreno}{\textbf{rocha}}, usando radar que
              atravessa o gelo: para ela, a montanha termina na pedra.
        \item O Nepal mediu até a \textcolor{corAgua}{\textbf{neve}}: para ele, o cume é
              onde o alpinista pisa.
      \end{itemize}
      \medskip
      Essa parte da briga era de \alert{definição}, não de medida.
    \end{column}
    \begin{column}{0.48\textwidth}
      \centering
      \begin{tikzpicture}[scale=0.95]
        \fill[corTerreno!30] (0.6,0) -- (2.9,2.55) -- (5.2,0) -- cycle;
        \draw[terreno] (0.6,0) -- (2.9,2.55) -- (5.2,0);
        \fill[corAgua!30] (2.35,2.05) .. controls (2.6,3.0) and (3.2,3.0) ..
                          (3.45,2.05) -- cycle;
        \draw[corAgua, thick] (2.35,2.05) .. controls (2.6,3.0) and (3.2,3.0) .. (3.45,2.05);
        \draw[corAgua, dashed] (2.35,2.05) -- (3.45,2.05);
        \draw[mRose, {Stealth[length=1.8mm]}-{Stealth[length=1.8mm]}] (3.75,2.05) -- (3.75,2.92);
        \node[rotulo, text=mRose, right] at (3.85,2.5) {a neve};
        \node[rotulo, text=corAgua, above] at (2.9,3.05) {cume de neve};
        \node[rotulo, text=corTerreno, right] at (3.5,1.5) {rocha};
      \end{tikzpicture}
    \end{column}
  \end{columns}
\end{frame}

\begin{frame}{Segunda discordância: a partir de onde se conta?}
  \begin{columns}[T]
    \begin{column}{0.46\textwidth}
      Esta é a parte que interessa a esta aula. Cada país carregou seu zero por terra até
      o Himalaia:
      \medskip
      \begin{itemize}
        \item a \textbf{China} partiu do \bnum{Mar Amarelo}, a leste;
        \item o \textbf{Nepal} partiu da \bnum{Baía de Bengala}, ao sul.
      \end{itemize}
      \medskip
      Dois marégrafos diferentes, a milhares de quilômetros um do outro, cada um com seu
      nível médio do mar. \alert{Dois zeros, duas respostas} --- para o mesmo cume.
    \end{column}
    \begin{column}{0.50\textwidth}
      \centering
      \begin{tikzpicture}[scale=0.95]
        \fill[corTerreno!22] (0,0) rectangle (6.3,1.5);
        \fill[corTerreno!35] (2.45,1.5) -- (3.15,2.85) -- (3.85,1.5) -- cycle;
        \draw[terreno] (2.45,1.5) -- (3.15,2.85) -- (3.85,1.5);
        \fill[mRose] (3.15,2.85) circle (0.075);
        \node[rotulo, above] at (3.15,2.95) {Everest};
        \draw[corAgua, very thick] (0,0.55) -- (0.95,0.55);
        \node[rotulo, text=corAgua, below, align=center] at (0.5,0.45) {Baía de\\Bengala};
        \draw[corAgua, very thick] (5.35,0.9) -- (6.3,0.9);
        \node[rotulo, text=corAgua, below, align=center] at (5.85,0.8) {Mar\\Amarelo};
        \draw[corAgua, dashed] (0.95,0.55) -- (3.15,0.55);
        \draw[corAgua, dashed] (5.35,0.9) -- (3.15,0.9);
        \draw[mRose, {Stealth[length=1.8mm]}-{Stealth[length=1.8mm]}] (3.15,0.55) -- (3.15,0.9);
        \node[rotulo, text=mRose, left] at (3.0,0.73) {zeros diferentes};
      \end{tikzpicture}
    \end{column}
  \end{columns}
\end{frame}

\begin{frame}{E ainda: a montanha entorta o próprio zero}
  \begin{columns}[T]
    \begin{column}{0.50\textwidth}
      O Himalaia é uma quantidade enorme de massa empilhada. Essa massa \bnum{puxa o fio
      de prumo} e \bnum{levanta o geoide} embaixo dela.
      \medskip

      Ou seja: a superfície de referência sob o Everest \alert{não está onde um modelo
      global simples colocaria}.
      \medskip

      Sem medir a gravidade \emph{ali}, a altura sai errada --- por metros, não por
      centímetros.
    \end{column}
    \begin{column}{0.46\textwidth}
      \centering
      \begin{tikzpicture}[scale=0.88]
        \fill[corTerreno!35] (1.6,1.1) -- (3.0,2.9) -- (4.4,1.1) -- cycle;
        \draw[terreno] (1.6,1.1) -- (3.0,2.9) -- (4.4,1.1);
        \fill[corTerreno!18] (0,0.6) rectangle (6.0,1.1);
        \draw[geoide] (0,0.35) .. controls (1.8,0.4) and (2.2,0.95) .. (3.0,0.95)
                      .. controls (3.8,0.95) and (4.2,0.4) .. (6.0,0.35);
        \draw[corElip, dashed] (0,-0.75) -- (6.0,-0.75);
        \node[rotulo, text=corElip, below] at (4.9,-0.8) {elipsoide};
        \node[rotulo, text=corGeoide, below, align=center] at (1.05,0.20)
              {geoide levantado\\pela montanha};
      \end{tikzpicture}
    \end{column}
  \end{columns}
\end{frame}

\begin{frame}{A campanha conjunta}
  \centering
  \small
  \renewcommand{\arraystretch}{1.4}
  \begin{tabular}{llp{6.0cm}}
    \toprule
    \textbf{22/05/2019} & Nepal & equipe liderada por Khim Lal Gautam chega ao cume com
      receptor GNSS e radar de penetração, e mede a espessura da neve \\
    \textbf{maio/2020} & China & equipe chinesa ao cume, com GNSS (incluindo BeiDou),
      radar e teodolitos \\
    \textbf{2019--2020} & ambos & nivelamento de precisão, levantamento gravimétrico e
      cálculo de um \bnum{geoide gravimétrico local} \\
    \textbf{08/12/2020} & juntos & anúncio conjunto do valor único \\
    \bottomrule
  \end{tabular}
  \bigskip

  Os dois países processaram \bnum{os dados juntos}, em vez de cada um publicar o seu.
\end{frame}

\begin{frame}{O número que encerrou a discussão}
  \begin{columns}[c]
    \begin{column}{0.46\textwidth}
      \begin{ideia}[title={Everest --- valor oficial desde 08/12/2020}]
        \centering
        {\Large \bnum{$8\,848{,}86$ m}}
        \medskip

        altura do cume de \textbf{neve}, acima de um \textbf{geoide gravimétrico}
        calculado \alert{com base no IHRS}
      \end{ideia}
    \end{column}
    \begin{column}{0.50\textwidth}
      \centering
      \small
      \renewcommand{\arraystretch}{1.25}
      \begin{tabular}{lr}
        \toprule
        Survey of India, 1954 (neve) & $8\,848$\,m \\
        China, 2005 (rocha) & $8\,844{,}43$\,m \\
        \midrule
        \textbf{conjunto, 2020 (neve)} & \bnum{$8\,848{,}86$\,m} \\
        \bottomrule
      \end{tabular}
      \medskip

      {\small A diferença entre o valor de 2020 e o chinês de 2005 --- $4{,}43$\,m ---
      é, em boa parte, a \textcolor{corAgua}{calota de neve} que um media e o outro não.}
    \end{column}
  \end{columns}
\end{frame}

\begin{frame}{O que mudou não foi a montanha}
  \begin{columns}[T]
    \begin{column}{0.48\textwidth}
      \begin{armadilha}[title={O que NÃO aconteceu}]
        O Everest não cresceu $86$ centímetros em 2020.
        \medskip

        (Ele cresce, sim --- alguns milímetros por ano, porque a Índia continua empurrando
        a Ásia. Mas não foi isso.)
      \end{armadilha}
    \end{column}
    \begin{column}{0.48\textwidth}
      \begin{ideia}[title={O que aconteceu}]
        Os dois países passaram a contar \bnum{a partir do mesmo zero} --- um zero que não
        é o porto de nenhum dos dois, mas o valor de potencial combinado pelo mundo
        inteiro.
        \medskip

        E mediram a gravidade local para saber onde esse zero passa, ali, embaixo da
        montanha.
      \end{ideia}
    \end{column}
  \end{columns}
  \bigskip
  \centering
  Uma disputa de setenta anos terminou quando os dois lados \bnum{concordaram sobre o
  zero}, não sobre o cume.
\end{frame}

% ==================================================================
\section{Fechamento}

\begin{frame}{Para levar para casa}
  \footnotesize
  \begin{enumerate}
    \setlength{\itemsep}{3.5pt}
    \item Altura só existe \bnum{a partir de alguma coisa}. Sem dizer de onde se conta,
          a pergunta ``qual é a altura?'' não tem resposta.
    \item O \textcolor{corElip}{\textbf{elipsoide}} é geometria pura --- é o que o GPS usa,
          e ele \alert{não sabe para onde a água corre}.
    \item O \textcolor{corGeoide}{\textbf{geoide}} é a superfície de nível que o mar
          escolheria. É o zero \emph{físico}, e vale a regra
          $\hgnss = \Horto + \Nond$.
    \item As superfícies de nível \alert{não são paralelas}: ``subir 1 metro'' não é
          sempre a mesma coisa. O que se salva é o \bnum{número geopotencial}
          $C = W_0 - W$, o esforço, que não depende do caminho.
    \item O \bnum{IHRS} escolheu um zero para o planeta ---
          $W_0 = 62\,636\,853{,}4\ \text{m}^2/\text{s}^2$ --- e o \bnum{IHRF} é a rede que
          o materializa.
    \item O Everest ficou em \bnum{$8\,848{,}86$\,m} não porque alguém mediu melhor, mas
          porque China e Nepal passaram a contar do mesmo lugar.
  \end{enumerate}
  \smallskip
  \begin{ideia}[title={Experimente}]
    \small No simulador \textbf{Modelos Terrestres}: abas \textbf{Geóide} e
    \textbf{Teluróide}.
  \end{ideia}
\end{frame}

% ==================================================================
\appendix

\begin{frame}[standout]
  Apêndice\\[6pt]
  {\normalsize Para quem quer ir além}
\end{frame}

\begin{frame}{A1 · O campo de gravidade}
  \small
  O potencial de gravidade $W$ soma a atração das massas e o efeito da rotação:
  \[ W = V + \Phi, \qquad \Phi = \tfrac{1}{2}\,\omega^2 (x^2 + y^2) \]
  \begin{itemize}
    \item $V$: potencial gravitacional (newtoniano) de toda a massa da Terra
    \item $\Phi$: potencial centrífugo, que é o que achata o planeta
    \item o vetor gravidade é $\vec{g} = \nabla W$, e $|\vec g| \approx 9{,}78$ a
          $9{,}83\ \text{m/s}^2$ do Equador ao polo
  \end{itemize}
  \medskip
  Uma \textbf{superfície de nível} é o lugar onde $W = \text{const}$; $\vec g$ é sempre
  perpendicular a ela. O \textbf{geoide} é a superfície $W = W_0$.
  \medskip
  A distância entre duas superfícies de nível vizinhas é $\mathrm{d}n = -\mathrm{d}W / g$:
  onde $g$ é maior, elas ficam mais juntas. É a não-paralelidade da Parte 1, em uma linha.
\end{frame}

\begin{frame}{A2 · Número geopotencial e os tipos de altitude}
  \small
  \[ C_P = W_0 - W_P = \int_{\text{geoide}}^{P} g \; \mathrm{d}n \]
  A integral é a mesma por qualquer caminho: $C$ é uma função de ponto, e é isso que faz o
  nivelamento fechar.
  \medskip
  \centering
  \renewcommand{\arraystretch}{1.6}
  \begin{tabular}{lll}
    \toprule
    altitude & fórmula & divisor \\
    \midrule
    ortométrica & $H = C / \bar{g}$ & gravidade real média ao longo do prumo \\
    normal & $H^{*} = C / \bar{\gamma}$ & gravidade normal média (elipsoide) \\
    dinâmica & $H^{\mathrm{dyn}} = C / \gamma_{45}$ & constante, a normal em $45^{\circ}$ \\
    \bottomrule
  \end{tabular}
  \flushleft
  \medskip
  $\bar g$ não é observável sem hipótese sobre a densidade das rochas --- é por isso que a
  altitude \emph{normal} (Molodensky) é preferida por muitos países: $\bar\gamma$ se
  calcula exatamente.
\end{frame}

\begin{frame}{A3 · As duas relações de conversão}
  \small
  \begin{columns}[T]
    \begin{column}{0.48\textwidth}
      \[ h = H + N \]
      \textbf{Stokes} --- $H$ ortométrica, $N$ ondulação geoidal, a superfície de
      referência é o \textbf{geoide}.
    \end{column}
    \begin{column}{0.48\textwidth}
      \[ h = H^{*} + \zeta \]
      \textbf{Molodensky} --- $H^{*}$ normal, $\zeta$ anomalia de altura, a superfície de
      referência é o \textbf{quase-geoide}.
    \end{column}
  \end{columns}
  \medskip
  $\zeta \approx N$ no mar (coincidem exatamente) e se afastam nas montanhas, onde a
  diferença chega à ordem de decímetros a poucos metros.
  \medskip
  \begin{ideia}[title={No simulador}]
    A aba \textbf{Teluróide} mostra a construção de Molodensky: a cada ponto $P$ da
    superfície física corresponde um $Q$ no teluróide com potencial normal igual ao
    potencial real em $P$, e $\overline{PQ} = \zeta$.
  \end{ideia}
\end{frame}

\begin{frame}{A4 · O que a Resolução nº 1 da IAG (2015) estabelece}
  \small
  \begin{enumerate}
    \item As coordenadas verticais são \textbf{números geopotenciais}
          $C_P = W_0 - W(P)$, referidos a uma superfície equipotencial do campo de
          gravidade terrestre.
    \item O valor convencional adotado é
          $W_0 = 62\,636\,853{,}4\ \text{m}^2\text{s}^{-2}$, substituindo o valor de 1998
          ($62\,636\,856{,}0 \pm 0{,}5$), incompatível com os modelos mais recentes.
    \item A referência espacial do ponto $P$ é dada por coordenadas $X$ do \textbf{ITRF}.
    \item A realização é o \textbf{IHRF}: rede global com densificações regionais e
          nacionais, na mesma estrutura do ITRF.
    \item A exatidão pretendida para os números geopotenciais globais é
          $\pm 1 \times 10^{-2}\ \text{m}^2\text{s}^{-2}$, cerca de $\pm 1$\,mm em
          altitude física.
  \end{enumerate}
\end{frame}

\begin{frame}{A5 · Fontes}
  \scriptsize
  \begin{itemize}
    \item \textbf{IHRS/IHRF} --- IAG Resolution No.~1, IUGG General Assembly 2015 (Praga).
          Sánchez et al., \emph{Towards a first realization of the International Height
          Reference System}; GGOS, \emph{IHRS Realization}; Sánchez \& Sideris,
          \emph{Vertical datum unification for the IHRS}, Geophys.~J.~Int.~209(2).
    \item \textbf{Altitudes dinâmicas nos Grandes Lagos} --- NOAA/NGS, \emph{North American
          Vertical Datum and International Great Lakes Datum}; NOAA, \emph{Dynamic Heights
          for the International Great Lakes Datum of 2020}. O IGLD 1985 usa altitudes
          dinâmicas; o NAVD 88 usa ortométricas de Helmert.
    \item \textbf{Everest} --- $8\,848$\,m: Survey of India, 1954 (neve).
          $8\,844{,}43$\,m: China, 2005 (rocha, com radar de penetração no gelo).
          Cume nepalês em 22/05/2019, equipe de Khim Lal Gautam; equipe chinesa em maio de
          2020; anúncio conjunto em 08/12/2020 de $8\,848{,}86$\,m, referido a geoide
          gravimétrico calculado com base no IHRS.
    \item \textbf{Ondulações extremas} --- as mesmas citadas pelo simulador,
          \texttt{models/geoide.js}, a partir de EGM2008.
  \end{itemize}
  \medskip
  \normalsize
  \centering
  Os fatos e números desta aula foram conferidos em 21/09/2026.
\end{frame}

\end{document}
